\section{Background}
\label{Sec:background}
In this section, we present the preliminaries of distributed constraint optimization problems (DCOPs) and the version of belief propagation used to solve them -- the min-sum algorithm and its version that includes damping.
\subsection{Distributed Constraint Optimization Problems (DCOPs)}
\label{Sec:DCOP_defs}
\sloppy
A DCOP\footnote{For lack of space and to avoid redundancy, we do not provide a separate description of a centralized COP. A COP can be considered a DCOP with a single agent.} is a tuple $\langle{\mathcal A},{\mathcal X},{\mathcal D},{\mathcal C}\rangle$, where ${\mathcal A}$ is a finite set of agents $\{ A_1,A_2,\ldots,A_n \}$; ${\mathcal X}$ is a finite set of variables $\{ X_1,X_2,\ldots,X_m \}$, where each variable is held by a single agent (an agent may hold more than one variable); ${\mathcal D}$ is a set of domains $\{ D_1, D_2,\ldots,D_m \}$, where each domain $D_i$ contains the finite set of values that can be assigned to variable $X_i$ and we denote an assignment of value $d \in D_i$ to $X_i$ by an ordered pair $\langle X_i, d \rangle$; and ${\mathcal C}$ is a set of constraints (relations), where each constraint $C_j \in {\mathcal C}$ defines a non-negative {\em cost} for every possible value combination of a set of variables and is of the form $C_j : D_{j_1} \times D_{j_2} \times \ldots \times D_{j_k} \rightarrow \mathbb{R}^+ \cup \{0\}$. A {\em binary constraint} refers to exactly two variables and is of the form $C_{ij} : D_i \times D_j \rightarrow \mathbb{R}^+ \cup \{0\}$.
In our discussion of Min-sum we often refer to $C_{ij}$ as the {\em cost table} and to $C_{ij}[w,r]$ as the entry corresponding to $w \in D_i$ and $r \in D_j$.

A {\em binary DCOP} is a DCOP in which all constraints are binary. Agents are {\em neighbors} if they are involved in the same constraint. A {\em partial assignment} (PA) is a set of value assignments to variables, in which each variable appears at most once. A constraint $C_j \in {\mathcal C}$ of the form $C_j : D_{j_1} \times D_{j_2} \times \ldots \times D_{j_k} \rightarrow \mathbb{R}^+ \cup \{0\}$ is {\em applicable} to $PA$ if each of the variables $X_{j_1} , \ldots , X_{j_k}$ is included in PA. The {\em cost of a partial assignment} $PA$ is the sum of all applicable constraints to $PA$ over the value assignments in $PA$. A {\em complete assignment} (i.e., {\em solution}) is a partial assignment that includes all variables ($\mathcal X$). An {\em optimal solution} is a complete assignment with minimal cost.
For simplicity, we assume that each agent holds exactly one variable (i.e., $n=m$) and focus on binary DCOPs. These assumptions are common in DCOP literature (e.g.,~\cite{Petcu05b,YeohFK10}).
Min-sum's complexity is exponential in the arity of constraints, and therefore its applicability to problems that include constraints with high arity is limited~\cite{FarinelliRPJ08,YedidsionZF18}. However, this limitation is not relaxed by the use of damping, therefore, it is out of the scope of this study. We note that in contrast to Attentive-DMS, the versions of Min-sum we propose in this study do apply to factor-graphs with non binary function-nodes.
\subsection{Min-sum Belief Propagation}
\label{secmaxsum}
Min-sum operates on a \emph{factor-graph}, a bipartite graph in which the nodes represent variables and functions (constraints)~\cite{KschischangFL01}.\footnote{We preserve the terminology of~\cite{FarinelliRPJ08} and call constraint representing nodes in the factor-graph ``function-nodes".}
Each variable-node, representing a variable of the original DCOP, is connected to all function-nodes representing constraints that it is involved in. In other words, a function-node is connected to all variable-nodes involved in the constraint it represents.
Variable-nodes and function-nodes take an active role in Min-sum, i.e., they can send and receive messages, and can perform computation. When used to solve a DCOP, each node's role is performed by an autonomous agent.
In Min-sum, a message sent to -- or from -- variable-node $X$ (for simplicity, we use the same notation for a variable and the variable-node representing it) is a vector of size $|D_X|$ including a cost (or the belief of a cost) for each value in $D_X$.
In the first iteration, all nodes assume that all messages they previously received (in iteration $0$) include vectors of zeros. Each iteration is \emph{two-phase}: first all variable-nodes compute and send their messages, then all function-nodes compute and send theirs using the variable messages of the same iteration (this is the schedule implemented and analyzed throughout the paper). A message $Q_{X \rightarrow F}^i$ that variable-node $X$ sends to function-node $F$ in iteration $i$ is formalized as follows:
$Q_{X \rightarrow F}^i =\underset{F'\in F_X \setminus \{ F\} }{ \sum} R_{F' \rightarrow X}^{i-1} - \alpha$,
where $F_X$ is the set of function-node neighbors of variable-node $X$ and $R_{F' \rightarrow X}^{i-1}$ is the message sent to variable-node $X$ by function-node $F'$ in iteration $i-1$.
$\alpha$ is a constant that is reduced from all costs included in the message in order to prevent the costs from growing arbitrarily large.
A message $R_{F \rightarrow X}^{i}$ sent from a function-node $F$ to a variable-node $X$ in iteration $i$,
includes for each value $x \in D_X$:
$R_{F \rightarrow X}^{i} = \underset{{PA_{-X}}} \min\, cost (\langle X , x \rangle, PA_{-X})$,
where $PA_{-X}$ is a possible combination of value assignments to variables involved in $F$ not including $X$. $cost(\langle X , x \rangle, PA_{-X})$ represents the cost of a partial assignment $a = \{ \langle X, x \rangle, PA_{-X} \}$, which is:
$f(a) + \underset{X' \in X_F \setminus \{X\},\langle X',x' \rangle \in a}{\sum} (Q_{X' \rightarrow F}^{i})[x']$,
where $f(a)$ is the original cost in the constraint represented by $F$ for the partial assignment $a$, $X_F$ is the set of variable-node neighbors of $F$, and $(Q_{X' \rightarrow F}^{i})[x']$ is the cost in the message sent from variable-node $X'$ in the first phase of the same iteration, for the value $x'$ that is assigned to $X'$ in $a$.
$X$ selects its value assignment $\hat{x} \in D_X$ following iteration $k$ as follows:
$\hat{x} = \underset{x\in D_X}{\mathrm{argmin}} \sum_{F\in F_X} (R_{F \rightarrow X}^k)[x]$.
\subsection{Damped Min-sum (DMS)}
In order to add damping to Min-sum, a parameter $\lambda \in [0,1)$ is used. Before sending a message in iteration $k$, a node performs calculations as in standard Min-sum. We use $\widehat{m^k_{i \rightarrow j}}$ to denote the result of the calculation made in standard Min-sum for the content of a message intended to be sent from node $i$ to node $j$ in iteration $k$ and $m^{k-1}_{i \rightarrow j}$ to denote the message sent by $i$ to $j$ at iteration $k-1$. The message sent by $i$ to $j$ at iteration $k$ is calculated as follows:
$m^{k}_{i \rightarrow j} = \lambda m^{k-1}_{i \rightarrow j} + (1 - \lambda) \widehat{m^k_{i \rightarrow j}}$.
Thus, $\lambda$ expresses the weight given to previous calculations performed with respect to the most recent calculation performed. When $\lambda = 0$, the resulting algorithm is standard Min-sum.
\subsection{Backtrack Cost Trees}
\label{sec:bct}
	For analyzing the behavior of Min-sum, the use of a {\em backtrack cost tree} (BCT) was proposed by \citet{ZivanLG20}. It allows tracing for each belief the entries in the cost tables held by function-nodes that were used to compose it. In other words, the components of the assignment's cost.
A BCT is defined for a belief that appears in a message (either from some variable $X_{i}$ to a function-node $F_{ij}$, or from some function-node $F_{ij}$ to a variable $X_{i}$). The belief is on the cost of assigning some value $x\in D_{i}$ to variable $X_{i}$. Without loss of generality, we will elaborate on messages from variables to function-nodes.
The belief is a sum of various components from which the BCT is composed. The root is the cost of a decision to assign some value to a variable at a time $t$ and the directed edges from its children in the tree include the beliefs that were summed to generate that belief. These edges lead to nodes representing the neighbors from which the parent node received messages in time $t-1$. Each of those nodes is connected to the nodes from which it received messages at time $t-2$, with the edges containing the beliefs that were passed to it. The tree leaves represent the nodes at time $0$.
